\documentclass[11pt]{article}
\usepackage[T1]{fontenc}
\usepackage{lmodern,amsmath,amssymb,amsthm,booktabs}
\usepackage[margin=1in]{geometry}
\usepackage[hidelinks]{hyperref}
\newtheorem{proposition}{Proposition}
\title{Sharing auxiliary roles without losing the finite-network rank deficit}
\author{Douglas Colkitt\thanks{Prepared with assistance from OpenAI Codex. This is not independent review or endorsement by OpenAI.}}
\date{October 7, 2026}
\begin{document}
\maketitle
\begin{abstract}
We identify first-stage and third-stage side roles in the finite bit network
used by OpenAI's integer-multiplication manuscript. An explicit matching makes
the new frame transitions nested, removing one third of the side roles while
preserving the absolute rank deficit. Combined with our previous routing and
parameter audits, this supports the conditional exponent saving $\kappa=2^{-75}$.
The complex network is unchanged. The improvement concerns an asymptotic
exponent; no practical runtime comparison is claimed.
\end{abstract}

\section{Setting and the matching}
We retain the finite-alphabet, fixed finite number of one-dimensional tapes
model of \emph{Integer multiplication below $n\log n$}, pinned at
\begin{center}\texttt{adc7f1241b42e322a6451854ab7e4b4c146bf78a}.\end{center}
All unchanged algorithmic interfaces remain assumptions. We use the previously
audited $h=46$ instance. Let $\mathcal T$ be the triples of $[h]$,
$v=\binom h3$, $N=v^3$, $m=h^3$, and $z=3\binom{h-3}{2}$.
Two triples are neighbors when their intersection has size one.
The rational label space is $F=\mathbb Q^h$ with form $I-J/9$;
triple indicators have norm two and neighboring indicators are orthogonal.

Pair consecutive ground points. Define a permutation $\pi$ of $\mathcal T$.
If a triple contains a full pair and a singleton, keep the singleton and cycle
its full pair among the pairs other than the singleton's pair. Otherwise it
uses three different pairs: keep the point in the least indexed pair and flip
the other two points to their partners. The first rule is a cycle for each fixed
singleton; the second is an involution. Their classes are disjoint, and
\[
 |T\cap\pi(T)|=1\quad\text{for every }T\in\mathcal T.
\]
Thus $\pi$ is explicit and $t_T\perp t_{\pi(T)}$.

Fix a common ordering of triples, and let $n_S(j)$ be the $j$th neighbor of $S$.
A stage-1 side role has fixed second and third triples $(A,B)$ and physical
first-factor target/source triples $(S,n_S(j))$. Pair it with the stage-3 side
role with fixed first and second triples $(S,\pi(A))$ and third-factor
target/source triples $(B,n_B(j))$. This is bijective: recover $A$ with
$\pi^{-1}$ and $j$ from the final neighbor list. There are $R=Nz$ such pairs.
Only these side roles are shared; stage-2 side roles and all central roles
remain separate.

\section{Scalar correctness and frame accounting}
Every first-stage invocation finishes before the third stage starts. A scalar
invocation restores all its auxiliary values, regardless of their initial
values. Identifying the paired roles therefore preserves the scalar bank
exchange and restoration of every remaining auxiliary value. No simultaneous
roles within an invocation are identified.

Scalar restoration alone does not prove an improvement in the rank count.
The last stage-1 gate on the shared role has label
\[
 E=F\otimes\langle t_A\rangle\otimes\langle t_B\rangle,\qquad \dim E=h.
\]
The first stage-3 gate has label
\[
 H=\langle t_S\otimes t_{\pi(A)}\rangle^\perp_{F\otimes F}
       \otimes\langle t_B\rangle,\qquad \dim H=h^2-1.
\]
For any $x\in F$, the inner product of $x\otimes t_A$ with
$t_S\otimes t_{\pi(A)}$ vanishes. Hence $E\subset H$. Both labels are
nondegenerate under the source's tensor decompositions.

\begin{proposition}[Exact saving per shared role]
The identification removes one role and exactly $m$ units of total rational
edge rank, while preserving all gate compatibility and endpoint conditions.
\end{proposition}
\begin{proof}
The old terminal edges are $E\to F^{\otimes3}$ and $0\to H$. Their total
projection-difference rank is $(m-h)+(h^2-1)$. The new edge $E\to H$ has
rank $h^2-1-h$, exactly $m$ less. All other edges and gate labels remain
unchanged. The shared role still has initial matrix zero and final matrix
$I_m$ and is fixed by the completed scalar permutation. All data endpoints
are unchanged. Thus the identity
$M_{\mathrm{sink}(\rho(w))}-M_{\mathrm{source}(w)}=I_m$ holds on every role.
The common-frame and rational-shear arguments apply to this fixed circuit.
\end{proof}

Consequently
\begin{align*}
 R&=9475984020888000,\\
 W_{\rm b}^*&=W_{\rm b}-R=18958995769111200,\\
 s_{\rm b}^*&=s_{\rm b}-mR=1845392811609813681600,\\
 W_{\rm b}^*m-s_{\rm b}^*&=N-2L_{\rm b}>0.
\end{align*}
The relative deficit increases because the absolute deficit is preserved while
the denominator decreases. The resulting new relative deficit is
\[
 \eta_{\rm b}^*=\frac9{29015910268}.
\]
This changes a network hypothesis of the previous fixed-network tuning bound;
it does not contradict that bound.

\section{Conditional multiplication witness}
Put
\[
 a=\frac{27}{10^{12}},\qquad a_{\rm c}=\frac9{500000000000},
 \qquad L=\frac{5743}{500}.
\]
The existing exact logarithm certificate proves $\log m<L$, and exact arithmetic
proves $\eta_{\rm b}^*>aL$. Therefore
$m^{1-a}>m(1-aL)>m(1-\eta_{\rm b}^*)$.
The unchanged complex network supports $\sigma=1-a_{\rm c}$ as before.
Choose
\begin{gather*}
 \tau=1-a,\quad \sigma=1-a_{\rm c},\quad
 \epsilon=\frac1{21},\quad \beta=\frac9{10},\quad c=\frac{9a}{10},\\
 \lambda=1-\frac{19a^2}{20},\quad \lambda'=1-\frac{9a^2}{10},\quad
 \delta=\frac1{16},\quad C_1=20.
\end{gather*}
The packed condition is $\tau(1+c/\beta)=1-a^2<\lambda<\lambda'<1$.
The separate conditions involving $\sigma$ follow from
$19a^2/20<a_{\rm c}$ and $9a^2<a_{\rm c}$. The seven margins are
\[
 g_2=g_3=G=\frac{3a^2}{70}
 =\frac{2187}{70000000000000000000000000}>2^{-75},
\]
with $g_1>1/2$, $g_4=20a/21$, $g_5=13/112$, $g_6=299/336$ and
$g_7=1/21$, all larger than $G$.
For $\rho=G-2^{-75}>0$, the existing strict-gap argument absorbs fixed
logarithmic factors and yields $O(n(\log n)^{1-2^{-75}})$ conditionally.

The dimension exponent, stopping threshold and guard bound match the 2$^{-76}$
tuning note. Thus the guard is $p^{20/21}=o(p)$,
$\alpha=\Theta(p^{23/84})$ and $\gamma=O(p^{25/42})=o(p)$.
The enlarged chunk has $K=\Theta(p^{3a/70})$, still $K=o(\ell)$ and
$K/\log p\to\infty$, where $\ell=\Theta(p^{20/21})$.
Record length remains superpolynomial. The bit-network change affects only
finite constants in its fixed-tape implementation; the complex network, its
arithmetic guard analysis, the numerical maps and final rounding are retained.
No appeal is made to the separate recurrence refinement, which had assumed
$\sigma\le\tau$; here $\sigma>\tau$ and the original layer inequalities pass.

\section{Artifacts and review boundary}
The companion repository supplies an exact parameter and count certificate,
a standalone patch, and finite checks:
\begin{itemize}
\item \path{scripts/reuse_network.py} and \path{certificates/stage-reuse.json};
\item \path{patches/h46-shared-side-75.patch};
\item \path{tests/test_reuse_network.py}.
\end{itemize}
The matching is checked on every one of the 15,180 triples at $h=46$.
Small rational projection checks verify the orthogonality used in the nesting
argument. A complete $h=6$ scalar model checks the bank exchange and restoration
of 151,200 auxiliary roles with nonzero initial values. These
finite tests support the written proof; they do not independently establish the
upstream algorithm or its time bounds. No global optimality claim is made.
\end{document}
